\section*{Hoofdstuk \ref{ch:b3:heterocycles} --- Heterocyclische chemie}

\begin{solution}{exo:b3:heterocycles:1}
Oxazool: O geeft er twee (één vrij elektronenpaar in het $\pi$-systeem, één in het vlak), N één, de
drie koolstofatomen drie: 6. Thiazool: hetzelfde met S: 6. Pyrimidine:
één van elk van de zes atomen, beide vrije elektronenparen van stikstof
in het vlak: 6. Purine: negen atomen, waarbij het stikstofatoom N--H er
twee geeft en de andere elk één: 10, een hückelgetal
($4n + 2$, $n = 2$).
\end{solution}

\begin{solution}{exo:b3:heterocycles:2}
Piperidine (11.1) $>$ DMAP (9.6) $>$ imidazool (7.0) $>$ pyridine (5.2) $\gg$ pyrrool (ongeveer $-4$,
geprotoneerd op koolstof).
\end{solution}

\begin{solution}{exo:b3:heterocycles:3}
2-Broompyrrool (pyrrool is zo reactief dat het gemakkelijk
meervoudig gebromeerd wordt); 2-broomfuraan; 3-broomindool.
\end{solution}

\begin{solution}{exo:b3:heterocycles:4}
2,5-Dimethylfuraan; 1,2,5-trimethylpyrrool; 2,5-dimethylthiofeen.
\end{solution}

\begin{solution}{exo:b3:heterocycles:5}
2-Methoxypyridine. Additie van methoxide op C2 zet de negatieve lading
van het \omterm{def:b3:heterocycles:snar}{meisenheimercomplex} op stikstof; op C3 kan ze alleen over
koolstofatomen verdeeld worden, zodat het intermediair, en de
overgangstoestand ervoor, veel hoger in energie liggen.
\end{solution}

\begin{solution}{exo:b3:heterocycles:6}
2-Aminopyridine (als natriumzout, daarna geprotoneerd bij de opwerking).
Het amide-ion addeert op C2, wat een anionisch $\sigma$-adduct geeft met H en $\ce{NH2}$ op C2 en de negatieve
lading op het ringstikstofatoom; het verliest hydride, dat een proton
van de aminogroep opneemt en als $\ce{H2}$ vertrekt.
\end{solution}

\begin{solution}{exo:b3:heterocycles:7}
Het zuurstofatoom van het N-oxide geeft een vrij elektronenpaar terug aan
de ring: resonantiestructuren dragen een negatieve lading op C2, C4 en
C6. Elektrofielen vallen C4 aan (C2 ligt dicht bij het positief geladen
stikstofatoom). Het zuurstofatoom wordt daarna verwijderd met
fosfortrichloride, wat 4-nitropyridine geeft.
\end{solution}

\begin{solution}{exo:b3:heterocycles:8}
2-Hydroxypyridine (een OH op C2) en 2-pyridon (N--H en C=O). Het pyridon behoudt zes
$\pi$-elektronen in de ring: zijn zwitterionische resonantiestructuur, met $\mathrm{O^-}$ en een
$\mathrm{N^+}$ dat zijn vrije elektronenpaar deelt, is een pyridiniumolaat; het is een amide waarvan
het vrije elektronenpaar van stikstof het sextet vervolledigt, en de
sterke waterstofbruggen van C=O en N--H begunstigen het in water.
\end{solution}

\begin{solution}{exo:b3:heterocycles:9}
2-Nitrobenzaldehyde, twee equivalenten methylacetoacetaat, en ammoniak:
de calciumkanaalblokker nifedipine.
\end{solution}

\begin{solution}{exo:b3:heterocycles:10}
Het eenhydrazine naar de groep $\ce{CH2}$ (inwendig, sterker gesubstitueerd) geeft
2,3-dimethylindool; dat naar de methylgroep (eindstandig) geeft
2-ethylindool. Zwakke zuren geven vooral het eerste, via het stabielere
eenhydrazine; sterkere zuren verhogen het aandeel 2-ethylindool.
\end{solution}

\begin{solution}{exo:b3:heterocycles:11}
Bij $\mathrm S_\mathrm NAr$ is de snelheidsbepalende stap de additie, die de binding C--X
niet verbreekt; het elektronegatievere fluor stabiliseert de anionische
overgangstoestand en het \omterm{def:b3:heterocycles:snar}{meisenheimercomplex} sterker. Bij $\mathrm S_\mathrm N2$ breekt de binding C--X in de
overgangstoestand, en vertrekt de zwakke binding C--I het snelst.
\end{solution}

\begin{solution}{exo:b3:heterocycles:12}
Pyridine: drie signalen in de verhouding 2\,:\,1\,:\,2, het eerste rond \qty{8.6}{ppm}. Pyrrool: twee gelijke
signalen bij 6.7 en \qty{6.2}{ppm} en een brede N--H rond \qty{8.1}{ppm}. Furaan: twee gelijke signalen bij
7.4 en \qty{6.4}{ppm}, zonder N--H. Thiofeen: twee gelijke signalen bij 7.33 en \qty{7.12}{ppm}, dicht
bij elkaar, tussen die van furaan.
\end{solution}

\begin{solution}{pb:b3:heterocycles:1}
\textbf{1.} $\ce{C6H5NHNH2 + (CH3)2CO -> C6H5NHN=C(CH3)2 + H2O}$.

\textbf{2.} Het aminestikstofatoom addeert aan het carbonylkoolstofatoom,
waarna water verloren gaat, zoals bij de vorming van een imine; zuur
activeert de carbonylgroep en maakt van de OH van het adduct een
uittredende groep.

\textbf{3.} De eindstandige $\ce{NH2}$: het vrije elektronenpaar van het andere
stikstofatoom is geconjugeerd met de fenylring en sterker gehinderd.

\textbf{4.} $\qty{108.0}{g/mol}$; $\qty{10.0}{g}/\qty{108.0}{g/mol} = \qty{0.0926}{mol}$.

\textbf{5.} Het hydrazon $\ce{C9H12N2}$, \qty{148.0}{g/mol}: $\qty{0.0926}{mol} \times \qty{148.0}{g/mol} = \qty{13.7}{g}$.

\textbf{6.} $\ce{C6H5-NH-NH-C(CH3)=CH2}$.

\textbf{7.} Het ringkoolstofatoom \textit{ortho}, het koolstofatoom \textit{ipso}, de twee stikstofatomen, het
hydrazonkoolstofatoom en de eindstandige $\ce{CH2}$.

\textbf{8.} De binding N--N breekt; een binding C--C vormt zich tussen het koolstofatoom \textit{ortho} en de $\ce{CH2}$.

\textbf{9.} Zes elektronen, alle componenten \omterm{def:b3:pericyclic:faciality}{suprafaciaal}: thermisch toegestaan.

\textbf{10.} Als lewiszuur katalyseert het de tautomerisatie, bindt het
een stikstofatoom en verzwakt het de binding N--N, en helpt het
ammoniak vertrekken.

\textbf{11.} Het koolstofatoom \textit{ortho} is $sp^3$ geworden, met een waterstofatoom: het verlies van dat proton
herstelt de benzeenring, een sterke drijvende kracht.

\textbf{12.} De gevormde aniline-$\ce{NH2}$ valt het iminekoolstofatoom aan en sluit de
vijfring (een 2-amino-indoline).

\textbf{13.} Ammoniak; de vorming van de aromatische pyrroolring van indool
drijft de stap.

\textbf{14.} $\ce{C9H12N2 -> C9H9N + NH3}$.

\textbf{15.} \qty{131.0}{g/mol}.

\textbf{16.} $\qty{0.0926}{mol} \times \qty{131.0}{g/mol} = \qty{12.1}{g}$.

\textbf{17.} Op C3, de meest nucleofiele positie van indool, die hier vrij
is; aanval daar houdt de benzeenring in het intermediair intact.

\textbf{18.} $\ce{C6H5NHNH-C(=CH2)-CH2CH3}$ (eindstandig) en $\ce{C6H5NHNH-C(CH3)=CH-CH3}$ (inwendig).

\textbf{19.} 2-Ethylindool en 2,3-dimethylindool.

\textbf{20.} Het inwendige is sterker gesubstitueerd; sterke zuren verhogen
het aandeel van het eindstandige.

\textbf{21.} 2,3-Dimethylindool, via het stabielere, sterker gesubstitueerde
eenhydrazine.

\textbf{22.} Acetaldehyde condenseert met zichzelf en de route via het
eenhydrazine geeft lage rendementen; indool wordt in plaats daarvan
gemaakt uit het hydrazon van pyrodruivenzuur, wat
indool-2-carbonzuur geeft, dat gedecarboxyleerd wordt.

\textbf{23.} Een mengsel van regio-isomeren moet gescheiden worden, wat
rendement kost; een symmetrisch keton of een route die de regiochemie
vastlegt, verdient de voorkeur.

\textbf{24.} 10.0 g fenylhydrazine kan hoogstens \qty{12.1}{g} 2-methylindool geven.
\end{solution}
